%%==========================================================================
\section{Introduction}
\label{sec:intro}
%%==========================================================================

The airline planning process is traditionally decomposed into a sequence of
sub-problems: timetable design, fleet assignment, crew rostering, tail
assignment, and recovery planning~\citep{barnhart2003applications,
gopalan1998mathematical}.  The \emph{tail assignment problem} (TAP)—assigning
a specific aircraft (identified by its tail number) to each flight leg—is the
penultimate planning step and is tightly coupled with aircraft maintenance
scheduling.  Poor tail assignment decisions directly inflate operational costs
through unnecessary ferry (repositioning) flights, suboptimal maintenance
routing, and safety-margin violations.

In addition to the basic assignment and continuity requirements, the TAP must
respect aircraft maintenance constraints.  The most common maintenance type is \emph{type-A}, which requires a maintenance check after a maximum flying-time interval.  Type-A maintenance is typically performed at a limited set of
airports, and the TAP must ensure that each aircraft is routed to a maintenance-capable airport before its flying-time limit is exceeded.  The type-A maintenance requirement is a key operational constraint in the TAP and has been studied extensively in the literature.  

\citet{khaled2018compact} proposed a compact 0--1 linear programming
formulation for the TAP without maintenance, featuring only $O(n_f n_p)$ binary variables (where
$n_f$ is the number of flights and $n_p$ the fleet size).  This is a
significant reduction over the classic multi-commodity flow model of
\citet{levin1971scheduling}, which requires $O(n_f^2 n_p)$ variables.
Further, \citet{khaled2018compact} extended the model to include type-A maintenance
requirements, which are defined by a maximum flying-time interval between maintenance checks, that adds other $O(n_f n_p D)$ 
binary variables, where $D$ is the number of calendar days in the planning horizon. 
Computational experiments reported in that paper show that instances with up
to 1\,500 flight legs and 40 aircraft can be solved to certified optimality
within one hour on a standard workstation.

\subsection*{Our contributions}

The present paper makes three contributions that check and correct
those results:

\begin{enumerate}
\item \textbf{Basic feasibility correction (C2--C4 / constraints (3)--(6)).}
  We show that the strict-time balance definition used in the compact basic
  model admits a simultaneous-departure corner case: one aircraft can be
  assigned to two flights departing at the same airport and same time without
  violating the balance inequalities.  We formalize this counter-example and
  enforce physical executability by adding explicit pairwise non-overlap
  constraints. 
  
  Our non-overlap correction retains the $O(n_fn_p)$ binary-variable space but
can add $O(n_f^2n_p)$ conflict inequalities in the worst case. In practical
timetables, only temporally conflicting pairs are instantiated; the resulting
constraint count is therefore governed by the sparse conflict set
$\mathcal O$ defined in \Cref{sec:basic}.


\item \textbf{Constraint~(13) correction.}  We show that the cumulative-flight-
  time maintenance constraint (Constraint~13 in \citet{khaled2018compact})
  contains a logical flaw in opening and one-endpoint intervals, and in the handling of consecutive flights/maintences. 
  We also show  by counterexample that a natural two-row endpoint split does not fully repair
  the model. We therefore formulate an exact cumulative-state recursion that
  propagates and resets the flight-hour counter without day-pair ambiguity.

\item \textbf{Event-based reformulation.}  The day-indexed type-A model
  above still repeats a calendar-day index across every maintenance
  variable. We give a complete, day-free reformulation of the same type-A
  C1--C14 requirements, in which each flight's own timestamp supplies all
  calendar information and maintenance is represented as a virtual flight
  triggered by, and indexed on, the preceding real flight rather than by
  calendar day. This removes the day dimension from the maintenance model's
  size and lets the flight-hour counter propagate exactly along each
  aircraft's actual route, with no day-pair window and therefore no
  possibility of the split-form failure mode proved for
  Constraint~(13).
\end{enumerate}

\subsection*{Paper organisation}

\Cref{sec:literature} reviews related work.  \Cref{sec:basic} restates the
basic TAP model and proves a C2--C4 corner-case correction.  \Cref{sec:maintenance}
presents the maintenance-extended model, the Constraint~(13) correction, and
the induced interpretation updates for C8--C14 under the corrected basic model.
\Cref{sec:event_reformulation} gives the day-free event-based reformulation
of the type-A requirements and proves its correctness directly.
\Cref{sec:computational} reports the numeric Constraint~(13) counterexample
and the A-only computational study. \Cref{sec:conclusion} concludes.
